\SourceLeaf{25}
\par
Attractio in directione coordinatis $y$ parallela et opposita, prodit mutando in his formulis $a$ in $b$, unde patet, has duas vires aequivalere unicae, cuius directio axi $2C$ normalis est, cuiusque intensitas invenitur, si in formula modo tradita $a$ in distantiam puncti attracti ab hoc axe mutatur.

Attractio denique in directione coordinatis $z$ parallela et opposita i. e. ad aequatorem normali, fit in casu, ubi $B=A$,
\[
=\frac{4\pi cAA}{CC}\int\frac{tt\mathrm{d}t}{1-(1-\frac{AA}{CC})tt}
\]
unde eruitur, si $C<A$, ponendo ut supra $\frac{C}{A}=\cos\varphi$,
\[
=\frac{4\pi c\cos\varphi}{\sin\varphi^3}(\operatorname{tang}\varphi-\varphi)
\]
si vero $C>A$, prodit
\[
\frac{4\pi cAAC}{(CC-AA)^{\frac32}}\log\frac{C+\BareRoot{(CC-AA)}}{A}-\frac{4\pi cAA}{CC-AA}
\]
Tandem, si omnes tres $A,B,C$ sunt aequales, i. e. si corpus est sphaera, attractiones secundum tres directiones principales fiunt
\[
\tfrac43\pi a,\quad\tfrac43\pi b,\quad\tfrac43\pi c
\]
i. e. identicae cum iis, quas nucleus sphaericus, in cuius superficie punctum attractum iacet, exerceret, si ipsius massa in centro esset concentrata. Hinc etiam sponte sequitur, puncta \textit{externa} a sphaera perinde attrahi, ac si eius massa tota esset in centro, uti \textsc{Newton} primus docuerat.

\par\bigskip\Needspace{6\baselineskip}
\pdfbookmark[1]{Additamentum}{attractionis-additamentum}
\centerline{ADDITAMENTUM.}\medskip
\par
Postquam haecce iam perscripta essent, innotuit, indicante ill. \textsc{Laplace}, commentatio egregia cl. \textsc{Ivory} in \textit{Philosophical Transactions} ad A.~1809; ubi idem argumentum per methodum ab iis, quibus usi erant ill. \textsc{Laplace} et \textsc{Legendre}, prorsus diversam tractatur. Summa elegantia ille geometra attractionem puncti externi ad attractionem puncti interni reducere docuit, i. e. problematis partem, quae semper pro difficiliori habita est, ad faciliorem. Methodus autem, per quam
